\documentclass[11pt]{article}
\usepackage[a4paper,margin=0.8in]{geometry}
\usepackage[T1]{fontenc}
\usepackage{lmodern}
\usepackage{amsmath,amssymb,booktabs,tabularx}
\usepackage{xcolor,listings}
\usepackage[hidelinks]{hyperref}
\usepackage{enumitem}
\setlength{\parindent}{0pt}
\setlength{\parskip}{5pt}
\lstset{language=R,basicstyle=\ttfamily\footnotesize,breaklines=true,
  columns=fullflexible,keepspaces=true,showstringspaces=false,
  backgroundcolor=\color{gray!7},frame=single,rulecolor=\color{gray!25},
  aboveskip=5pt,belowskip=5pt,xleftmargin=4pt,xrightmargin=4pt}
\newcommand{\code}[1]{\texttt{#1}}
\begin{document}
\begin{center}
{\Large\bfseries Assignment 3: Part A}\\[4pt]
{\large Guide to the \texttt{bnlearn} Functions}
\end{center}

The examples assume that \code{dat} contains the four factor variables
\code{Y}, \code{X1}, \code{X2}, and \code{X3}, and that
\code{new\_customer} contains predictor values with the same factor levels.
Report the installed version using \code{packageVersion("bnlearn")}.
Slide references below refer to Lecture Set 4.

\section*{A1. DAGs and factorization \normalfont\small (slides 15--19, 23)}
A Bayesian network factorizes a joint distribution as
\[
 p(y_1,\ldots,y_I)=\prod_{i=1}^{I}p(y_i\mid\Pi_i),
\]
where $\Pi_i$ is the parent set of node $i$.
\code{model2network()} converts a model string into a DAG. Each bracket
specifies a node and its parents: \code{[Y]} denotes a root node, while
\code{[X2|Y:X1]} corresponds to the factor $p(x_2\mid y,x_1)$.
The five candidate networks are:
\begin{lstlisting}
M1 <- model2network("[Y][X1|Y][X2|Y:X1][X3|Y:X1:X2]")
M2 <- model2network("[Y][X1|Y][X2|Y][X3|Y]")
M3 <- model2network("[Y][X1|Y][X2|Y:X1][X3|Y]")
M4 <- model2network("[Y][X1|Y][X2|Y][X3|Y:X1]")
M5 <- model2network("[Y][X1|Y][X2|Y][X3|Y:X2]")
\end{lstlisting}
For example, M3 represents
$p(y)p(x_1\mid y)p(x_2\mid y,x_1)p(x_3\mid y)$.
\code{modelstring()} performs the reverse conversion, returning the
string representation of a supplied network.
\begin{lstlisting}
modelstring(M3)
\end{lstlisting}
\code{naive.bayes()} constructs a classifier whose class node is the parent
of every predictor. The argument \code{training = "Y"} identifies the
class variable; the data supply the variable names and categories.
\begin{lstlisting}
nb <- naive.bayes(dat, training = "Y")
\end{lstlisting}
Its DAG is the same as M2: the attributes are mutually independent given Y.
The returned object also has a specialized classifier class used by
\code{predict()}.

\code{nparams()} counts free parameters. For discrete networks, the count
is $\sum_i Q_i(C_i-1)$, where $C_i$ is the number of states of node $i$
and $Q_i$ is the number of parent configurations. The \code{data} argument
provides category information for an unfitted network.
\begin{lstlisting}
nparams(M2, data = dat)  # 2 + 3 + 3 + 3 = 11
\end{lstlisting}

\section*{A2. Parameter learning \normalfont\small (slides 29--33)}
\code{bn.fit()} estimates local distributions given a network and its data.
With \code{method = "mle"}, it returns
\[
 \widehat\theta_{ijk}=\frac{N_{ijk}}{N_{ij}},\qquad
 N_{ij}=\sum_k N_{ijk},
\]
where $N_{ijk}$ counts observations in state $k$ under parent configuration
$j$. For example, $\widehat P(X_1=\mathrm{high}\mid Y=A)=130/333=0.3904$.
\begin{lstlisting}
fit_mle <- bn.fit(M2, data = dat, method = "mle")
\end{lstlisting}

\newpage
With \code{method = "bayes"}, \code{iss} is the imaginary sample size,
corresponding to the lecture's global precision $\alpha$. The uniform BDe
prior assigns
\[
 \alpha_{ijk}=\frac{\alpha}{C_iQ_i},\qquad
 \alpha_{ij}=\sum_k\alpha_{ijk}=\frac{\alpha}{Q_i}.
\]
The posterior and its mean are
\begin{align*}
 \boldsymbol\theta_{ij}\mid D
 &\sim\operatorname{Dir}(\alpha_{ij1}+N_{ij1},\ldots,
                         \alpha_{ijC_i}+N_{ijC_i}),\\
 E(\theta_{ijk}\mid D)
 &=\frac{\alpha_{ijk}+N_{ijk}}{\alpha_{ij}+N_{ij}}.
\end{align*}
\begin{lstlisting}
fit_M2 <- bn.fit(M2, data = dat, method = "bayes", iss = 12)
\end{lstlisting}
For $X_1\mid Y=A$, the prior is $\operatorname{Dir}(2,2)$ and the observed
low/high counts are $(203,130)$. The posterior is
$\operatorname{Dir}(205,132)$, giving a posterior mean high-satisfaction
probability of $132/337=0.3917$.

\code{coef()} extracts fitted probability tables. For a Bayesian fit,
the entries are posterior means of $\theta_{ijk}$, not complete
Dirichlet distributions.
\begin{lstlisting}
coef(fit_M2$X1)
\end{lstlisting}

\section*{A3. Classification probabilities \normalfont\small (slides 24, 36--38)}
Classification requires $p(Y=j\mid x_1,x_2,x_3,D)$. Under naive Bayes,
these probabilities are obtained by normalizing
\[
 p(Y=j\mid D)\prod_{\ell=1}^{3}p(x_\ell\mid Y=j,D)
\]
over the possible classes.

\code{predict()} returns predicted classes. For an ordinary fitted network,
\code{node = "Y"} specifies the target, \code{data} supplies the attribute
evidence, and \code{prob = TRUE} attaches the classification probabilities
as the \code{"prob"} attribute. Its matrix has one row per class and one
column per observation.
\begin{lstlisting}
pred <- predict(fit_M2, node = "Y", data = new_customer,
                method = "exact", prob = TRUE)
attr(pred, "prob")
\end{lstlisting}

\begin{tabularx}{\linewidth}{@{}lX@{}}
\toprule
Method & Interpretation \\
\midrule
\code{parents} & Uses only the target node's parents. Since Y has no
parents, it returns $p(Y\mid D)$ and ignores the attributes. It is
unsuitable for the conditional classification required here. \\[4pt]
\code{bayes-lw} & Uses likelihood-weighting simulation to incorporate
attribute evidence. The argument \code{n} controls the simulation size;
finite simulations introduce Monte Carlo error. \\[4pt]
\code{exact} & Performs exact inference in the fitted discrete network
using junction trees and belief propagation, without Monte Carlo
approximation error. \\
\bottomrule
\end{tabularx}

\begin{lstlisting}
set.seed(20261002)
predict(fit_M2, node = "Y", data = new_customer,
        method = "bayes-lw", n = 100000, prob = TRUE)
\end{lstlisting}
Objects from \code{naive.bayes()} use a specialized prediction method.
We use ordinary networks from \code{model2network()} when comparing
the three methods above. Exact probabilities are deterministic, but
the package may randomly break ties between equally probable classes.

\newpage
\code{cpquery()} estimates the probability of an \code{event} conditional
on \code{evidence}. With \code{method = "ls"}, it simulates observations
and estimates the event frequency among those matching the evidence.
The argument \code{n} specifies the number of simulation draws.
\begin{lstlisting}
set.seed(20261002)
cpquery(fit_M2, event = (Y == "A"),
        evidence = (X1 == "high" & X2 == "low" & X3 == "high"),
        method = "ls", n = 100000)
\end{lstlisting}
Both \code{cpquery()} and \code{bayes-lw} exhibit Monte Carlo variability.
Resetting the seed makes a simulation reproducible; increasing its size
generally reduces approximation error. Probabilities obtained from
separate simulation calls need not sum to exactly one.

\section*{A4. Marginal likelihood and model comparison
\normalfont\small (slides 39--40)}
\code{score()} evaluates a network using a specified criterion.
With \code{type = "bde"} and \code{iss = 12}, it returns the log marginal
likelihood under the specified Dirichlet prior:
\[
 \log p(D\mid M)
 =\log\int p(D\mid\Theta,M)p(\Theta\mid M)\,d\Theta.
\]
\begin{lstlisting}
score(M2, data = dat, type = "bde", iss = 12)
\end{lstlisting}
The local likelihood structure and independent Dirichlet priors allow the
integral to factorize over nodes. Thus,
\[
 p(D\mid M)=\prod_i p(D_i\mid M),\qquad
 \log p(D\mid M)=\sum_i\log p(D_i\mid M).
\]
The argument \code{by.node = TRUE} returns these individual contributions,
whose sum is the whole-network score.
\begin{lstlisting}
score(M2, data = dat, type = "bde", iss = 12, by.node = TRUE)
\end{lstlisting}
With \code{type = "bic"}, the function instead returns
\[
 \mathrm{BIC}_{\texttt{bnlearn}}
 =\log p(D\mid\widehat\Theta,M)-\frac{k}{2}\log n,
\]
where $k$ is the number of free parameters and $n$ is the sample size.
Larger values are preferred in this convention.
\begin{lstlisting}
score(M2, data = dat, type = "bic")
\end{lstlisting}
BIC is a penalized likelihood criterion, not the exact Dirichlet log
marginal likelihood. We use BDe scores to calculate posterior model
probabilities. With equal model priors and $\ell_g=\log p(D\mid M_g)$,
\[
 P(M_g\mid D)=
 \frac{\exp(\ell_g-\ell_{\max})}
      {\sum_h\exp(\ell_h-\ell_{\max})}.
\]
Subtracting the maximum log score improves numerical stability without
changing the resulting probabilities.

\subsection*{References}
Lecture Set 4, slides 15--40; installed \code{bnlearn} help pages for
\code{model2network}, \code{modelstring}, \code{naive.bayes},
\code{nparams}, \code{bn.fit}, \code{coef}, \code{predict},
\code{cpquery}, and \code{score}. Online documentation:
\url{https://www.bnlearn.com/documentation/}.

\end{document}
